
\section{Related Work}
\subsection{Uplift modeling} \label{sec:uplift_model}
Existing uplift methods can be classified into three categories \cite{gutierrez2017causal}, i.e. the Two-Model methods \cite{radcliffe2007using,nassif2013uplift}, the Class-Transformation methods \cite{jaskowski2012uplift} and the methods that model uplift directly \cite{radcliffe2011real,rzepakowski2012decision,soltys2015ensemble}. 

The Two-Model methods construct two independent models for the two groups of data. One model infers the label using the data from the treatment group and the other model is for the control group data. However, \cite{radcliffe2011real} points out that the Two-Model methods may miss the uplift signal. Then The Class Transformation methods are introduced by \cite{jaskowski2012uplift,athey2015machine}, which aim to create a new target to approach the uplift. Then a single model is proposed to learn the new target. But Class Transformation methods usually require a balanced dataset between the control and treatment groups. The last type of uplift method aims to directly infer the uplift. The work \cite{zaniewicz2013support} proposes a method based on a modification of the SVM model and the work \cite{guelman2014optimal} focuses on k-nearest neighbors to do the uplift estimation. 

The aforementioned works assume that the data of the control group and treatment group are randomly collected. If the collecting data is naturally observed, besides uplift estimation, uplift modeling also needs to reduce the bias of the data from the treatment group and control group. The most popular methods of this type are the doubly robust learning methods \cite{emsley2008implementing,zhao2015doubly,jung2021estimating,chernozhukov2018double}. They usually adopt the Inverse Propensity Scoring (IPS) \cite{hirano2003efficient,kitagawa2018should} to re-weight each instance, aiming at making the uplift estimator unbiased. And some methods\cite{louizos2017causal,shalit2017estimating} which estimate individual treatment effect can be used to estimate the uplift.

However, aforementioned uplift methods assume that the uplift can be fully estimated by the individual features. As we have stated before, the social relationships between users are important for uplift estimation. Then Guo et. al.\cite{guo2020learning,guo2020counterfactual,ma2021deconfounding} first introduce the networked observational data to the problem of causal effects estimation. NetEst\cite{jiang2022estimating} formalize the networked causal effects estimation to a multi-task learning problem and HyperSCI\cite{ma2022learning} learning causal effects on hypergraphs. These methods prove that networked data is important for predicting causal effects. Although incorporating the graph data will violate the Stable Unit Treatment Value Assumption(SUTVA), following works \cite{sobel2006what} have pointed out that SUTVA is not plausible in real-world scenarios. We will follow these works \cite{guo2020learning,guo2020counterfactual} to incorporate the graph data but are distinct from them by addressing the label scarcity problem for graph-based uplift estimation.  

\subsection{Partial Label Learning}
Partial label learning deals with the problem that each sample is associated with a set of candidate labels, among which only one label is the ground-truth label to be predicted. Existing partial label learning methods can be classified into two categories: the averaged-based strategy and the identification-based strategy. The averaged-based strategy assumes that each candidate label contributes equally to the model training \cite{hullermeier2006learning,gong2017regularization}, which may suffer from the problem that the real label is overwhelmed by other labels. To overcome this drawback, the identification-based methods give different confidence to different candidate labels by learning the topological information \cite{zhang2015solving,zhang2016partial}. Existing partial-label learning is often applied to automatic face naming, object detection, and web mining. As far as we know, this is the first work to apply partial label learning to uplift modeling.

Other typical types of weakly-supervised learning  include incomplete label learning\cite{zhao2015semi-supervised} and inaccurate label learning\cite{wu2018a}. They are not closely related to our problem, which will not be discussed here. 

\subsection{Graph Neural Network}
Graph neural networks (GNNs), aiming to generalize neural networks to deal with graph data, have drawn increasing research interest recently \cite{bronstein2017geometric,zhang2020deep} and show effectiveness in various tasks \cite{wang2017community,fan2020graph,zhang2021graph}. Generally, current GNNs can be divided into two categories: spectral-based methods and spatial-based methods. Spectral-based GNNs are originated from signal processing and are commonly based on the Laplacian Matrix \cite{bruna2014spectral,defferrard2016convolutional,kipf2017semi}.  Spatial-based GNNs regard the graph convolution as the 'message-passing' framework in the spatial domain, i.e. defining the graph convolution as nodes aggregating information from neighborhoods  \cite{gilmer2017neural,hamilton2017inductive}. And \cite{feng2021should, sui2022causal, ding2022causal} have explored the causal inference problems with GNN-based models. More GNN-based models can be referred in recent surveys \cite{wu2020comprehensive,zhang2020deep}. However, only a few works focus on uplift estimation using GNNs and they do not address some specific and important problems when using GNNs on uplift estimation. 

